\section{Experiments}
\subsection{Datasets}

We used the BFT val dataset provided in the NetTrace paper as our dataset for this project. This dataset is a highly dynamic bird flock dataset, which comes from documentary,  containing 6 bird species, 25 videos, and a total of 5020 frames, covering various backgrounds such as oceans, mountains, and rivers. The dataset has been manually annotated with bounding boxes and traceIDs, and these annotations are stored in JSON files in the COCO dataset format.

\subsection{Settings}
\label{sec:exp_settings}

\paragraph{Data split and protocol.}
All experiments are conducted on the \textbf{BFT validation set} (25 videos, 5020 frames) following the split provided by NetTrack/BFT. We process each sequence independently and report the average tracking performance averaged over all sequences. Our OpenCV pipeline is \textbf{training-free} and runs in an \textbf{online} manner: at frame $t$ it only uses information from frames $\le t$.

\paragraph{Annotation and file format.}
BFT annotations are provided in COCO-style JSON with per-frame bounding boxes and \texttt{traceID}s. For evaluation, we convert both ground truth and predictions into \textbf{MOTChallenge} text format: \texttt{frame, id, x, y, w, h} (pixel coordinates). This unified interface is shared by our OpenCV pipeline and the NetTrack baseline, enabling identical evaluation scripts.

\paragraph{Hyperparameters.}
We use a fixed configuration for all sequences (no per-video manual tuning beyond Step-1 adaptive glare statistics). Key parameters for Steps 1-3 are summarized in Tab.~\ref{tab:opencv_settings_steps1_3}. To prevent the table from being truncated in a two-column layout, we present it as a \texttt{table*} spanning both columns.

\begin{table*}[t]
\centering
\small
\setlength{\tabcolsep}{4pt}
\renewcommand{\arraystretch}{1.05}
\begin{tabular}{p{0.20\textwidth}|p{0.77\textwidth}}
\hline
\textbf{Module} & \textbf{Main hyperparameters (used in all experiments)}\\
\hline
\textbf{Step-1: Preprocessing} &
Subtitle ROI: y0=0.74, y1=1.0; subtitle detection: $v_{\min}=150$, $\delta_{\text{th}}=8$, $s_{\text{low}}=55$; glare detection: base $v_{\min}=230$, base $\delta_{\text{th}}=24$, adaptive range $v_{\min}\in[210,245]$, $\delta_{\text{th}}\in[12,60]$; texture gating: $\text{edge\_density}_{\min}=0.020$; glare fallback: total ratio max=0.015, max component ratio=0.05; bilateral filter: $d=7$, $\sigma_c=35$, $\sigma_s=15$; LoG: $\sigma=1.2$, $k=3$ \\
\hline
\textbf{Step-2: Camera Motion Compensation} &
ROI mode: strips; strip fraction: 0.18; margin fraction: 0.02; max shift: 50px; moving threshold: 1.2px; phase correlation response thresholds: global 0.15, ROI 0.12; ECC fallback enabled (iterations=60, eps=1e-4); specular weight: 0.6; MAD outlier rejection: $k=3$, min inlier ROIs=3 \\
\hline
\textbf{Step-3: Candidate Generation} &
Normalized diff $\epsilon=3.0$; diff blur kernel: 7; MAD $k=8$; thresholds clamped [8,80]; specular weight: 0.45; hard suppression: spec ratio=0.35, fg ratio=0.45; LoG diff auto-enable: edge density thresh=0.06, spec ratio=0.05, fg ratio=0.18; KNN bg: history=300, dist2Threshold=400, learning rate=0.01 (still), 0.0 (high fg); bg-only suppression: support min=0.06, persistence min=1.6; morphology kernels: open=5, close=9, bridge dilate=5, bridge close=11; min area=80, max boxes=12, nested IoA thresh=0.92; persistence decay=0.85; global persistence filter: enabled, min=1.35 \\
\hline
\textbf{Step-4: Refinement} &
\texttt{min\_area\_ratio}=5e{-}5; \texttt{max\_box\_area\_ratio}=0.3; \texttt{min\_fg\_pixels}=20; \texttt{fg\_ratio\_min}=0.2; \texttt{fill\_ratio\_min}=0.10; \texttt{spec\_frac\_max}=0.6; \texttt{min\_side}=4; aspect ratio in $[0.10,10.0]$; \texttt{CENTER\_FRAC}=0.3 \\
\hline
\textbf{Step-5: Tracking} &
IoU gate $\ge 0.3$; center-dist gate $\le 1.5 \times$ diag; cost weights: $w_\text{iou}=0.7$, $w_\text{dist}=0.3$; \texttt{max\_age}=5, \texttt{min\_hits}=2; birth gate: IoU$<0.10$ and dist$>0.8\times$diag \\
\hline
\end{tabular}
\caption{Experimental hyperparameters of the OpenCV-based pipeline (fixed across all sequences). Steps 1-3 parameters are highlighted based on updated code.}
\label{tab:opencv_settings}
\end{table*}

\paragraph{Evaluation metrics.}
We report classic multi-object tracking metrics including \textbf{MOTA}, \textbf{IDF1}, \textbf{ID switches}, \textbf{false positives}, and \textbf{misses}. Associations are computed using an IoU-based matching threshold of 0.5 (i.e., matches require IoU$\ge 0.5$ in the evaluation distance matrix). For completeness, we also compute \textbf{HOTA} using TrackEval on the same MOT-format files when available, keeping the IoU threshold consistent.



\subsection{Implementation details}
\label{sec:exp_impl}

% (If not already in the preamble)
% \usepackage{algorithm}
% \usepackage{algorithmic}

\paragraph{Overview.}
Our OpenCV-based system is a five-stage, training-free pipeline:
(1) preprocessing (subtitle removal, glare suppression, feature extraction),
(2) camera motion compensation,
(3) motion-based candidate generation,
(4) ROI refinement with instance splitting and rule-based filtering,
and (5) online tracking with \textbf{Iterative Lucas--Kanade (IK / LK)} motion refinement and gated association.

\paragraph{Step-1: preprocessing.}
Given an input frame $\mathbf{I}^{\text{rgb}}_t$, we compute:
\begin{itemize}
    \item Smoothed grayscale $I_t$ via bilateral filtering
    \item LoG feature $L_t$ via Laplacian on Gaussian-blurred intensity
    \item Hard valid mask $V_t$ (subtitles as 0)
    \item Soft specular mask $S_t$ (glare as 0, birds as 255)
\end{itemize}
Masked regions are filled with Gaussian-blurred intensity ($\sigma=3$) to avoid artificial edges. Glare parameters are adapted using median statistics from first 10 frames.


\paragraph{Step-2: camera motion compensation.}
We estimate translation $\mathbf{T}_t$ from $(L_{t-1}, L_t)$ using three-stage cascade: (1) global phase correlation, (2) ROI consensus with MAD outlier rejection, (3) ECC fallback. The previous intensity $I_{t-1}$ is warped to $\hat{I}_{t-1\rightarrow t}$ using $\mathbf{T}_t$.

\paragraph{Step-3: motion-based candidate generation.}
We compute normalized difference $D^{\text{norm}}_t$ between $I_t$ and $\hat{I}_{t-1\rightarrow t}$, apply adaptive MAD threshold, optionally fuse with LoG difference and KNN background subtraction. Morphological cleanup (open/close/bridge) yields binary motion mask. Connected components become candidate boxes, filtered by geometry and scored by temporal persistence.


\paragraph{Step-4: candidate refinement (splitting + birdness filters).}
For each coarse ROI, we refine instances by:
(1) cleaning the foreground mask via open/close,
(2) obtaining confident centers by distance transform,
(3) applying marker-based watershed to split adjacent birds,
and (4) filtering each instance by area ratio, fill ratio, foreground occupancy, aspect ratio, minimum side length, and specular fraction.
This produces refined detections $\mathcal{B}_t$. Finally, we apply a global IoA-based suppression and relative-area filtering to remove residual nested or fragmented boxes within the same frame.

\paragraph{Step-5: online tracking with IK refinement + gated association.}
Each track stores a bounding box $b_{t-1}$ and a sparse set of keypoints $\mathcal{P}_{t-1}$ inside the track ROI.
At frame $t$, we first obtain a \emph{coarse} box prediction $\hat{b}_t$ using the last observed displacement (constant-displacement prior).
We then refine this prediction using \textbf{Iterative LK (IK)} within the predicted ROI.

\textbf{IK motion refinement.}
Let $R_{t-1}$ and $R_t$ be the cropped grayscale ROIs from $G_{t-1}$ and $G_t$ around $\hat{b}_t$.
We track keypoints $\mathcal{P}_{t-1}$ from $R_{t-1}$ to $R_t$ using pyramidal LK and remove outliers via a forward-backward consistency check.
From the remaining inliers, we compute a robust displacement $\Delta=(\Delta x,\Delta y)$ as the median flow.
The refined predicted box is obtained by shifting $\hat{b}_t$ by $\Delta$.
If the number of valid keypoints drops below a threshold (blur, occlusion, low texture), we re-initialize $\mathcal{P}_{t}$ with Shi--Tomasi corners in the current ROI.

\textbf{Association and track management.}
Given IK-refined predictions and detections $\mathcal{B}_t$, we compute a gated matching cost:
\[
\text{cost}(i,j)=w_\text{iou}(1-\text{IoU}_{ij})
+ w_\text{dist}\frac{\|p_i-p_j\|_2}{\text{diag}(b_i)},
\]
where $p$ is the box center and $\text{diag}(b_i)$ normalizes distances by track scale.
Pairs violating IoU and center-distance gates are discarded.
Matching is solved by a global greedy strategy over feasible pairs.
Unmatched tracks are kept for up to \texttt{max\_age} frames and removed afterward; tracks are output only after \texttt{min\_hits} confirmations.
Unmatched detections spawn new tracks subject to a birth gate to reduce duplicate IDs near existing tracks.

\begin{algorithm}[t]
\small
\caption{Online tracking with IK (Iterative LK) refinement}
\label{alg:ik_tracker}
\begin{algorithmic}[1]
\REQUIRE Grayscale frames $G_{t-1}, G_t$, detections $\mathcal{B}_t$, active tracks $\mathcal{T}_{t-1}$
\ENSURE Updated tracks $\mathcal{T}_t$ and MOT outputs at frame $t$

\FOR{each track $\tau \in \mathcal{T}_{t-1}$}
    \STATE Predict $\hat{b}_\tau \leftarrow b_\tau + \Delta^{last}_\tau$ \hfill (constant displacement prior)
    \STATE Crop ROIs $(R_{t-1},R_t)$ from $(G_{t-1},G_t)$ around $\hat{b}_\tau$
    \IF{$\mathcal{P}_\tau$ not initialized \OR $|\mathcal{P}_\tau|$ is small}
        \STATE Initialize $\mathcal{P}_\tau \leftarrow$ Shi--Tomasi corners in $R_{t-1}$
    \ENDIF
    \STATE Track $\mathcal{P}_\tau$ by pyramidal LK: $R_{t-1}\rightarrow R_t$
    \STATE Forward-backward check; keep inliers $\mathcal{P}_\tau^{in}$
    \IF{$|\mathcal{P}_\tau^{in}|$ sufficient}
        \STATE $\Delta_\tau \leftarrow \mathrm{median}(\mathbf{p}_t-\mathbf{p}_{t-1})$ over $\mathcal{P}_\tau^{in}$
        \STATE Refine prediction: $\hat{b}_\tau \leftarrow \hat{b}_\tau + \Delta_\tau$
        \STATE Store last displacement: $\Delta^{last}_\tau \leftarrow \Delta_\tau$
    \ENDIF
\ENDFOR

\STATE Build feasible pairs between refined predictions $\{\hat{b}_\tau\}$ and detections $\mathcal{B}_t$ using IoU and distance gates
\STATE Match feasible pairs by global greedy on cost (IoU + normalized center distance)
\STATE Update matched tracks; age unmatched tracks; spawn new tracks from unmatched detections (birth gate)
\STATE Output $(t, id, x, y, w, h)$ for confirmed tracks
\end{algorithmic}
\end{algorithm}

\paragraph{Outputs.}
For each video, predicted trajectories are exported in MOTChallenge format for evaluation under identical settings as the baseline.


\subsection{Evaluation}

The NetTrack baseline is evaluated following the official multi-stage inference pipeline released by the authors. All experiments use the pretrained weights provided with the framework. These weights are employed only for inference, and no additional training or fine-tuning is performed.

Specifically, the evaluation consists of two sequential stages. In the first stage, object detection is performed using Grounding DINO with the official pretrained model. The detector is applied to the BFT validation sequences using a text prompt corresponding to birds, and frame-level detection results are saved in text format. These detection outputs serve as the sole input to the subsequent tracking stage.

In the second stage, tracking is carried out using the official NetTrack demo script. The tracker loads the precomputed detection results and applies CoTracker, initialized with its pretrained weights, to propagate point-level motion across frames. Detection-to-track association is then performed by NetTrack to produce sequence-level tracking results with consistent identities. The final outputs are stored in MOTChallenge format.

After completing both detection and tracking stages, the resulting files are evaluated using the same evaluation script as for the proposed OpenCV-based pipeline. Metrics including IoU, HOTA, MOTA, and IDF1 are computed under identical settings, ensuring a fair comparison between the proposed model and the NetTrack baseline.